%------------------------------------------------------------------------------%
\begin{question}
  Calcular a integral \(\int_{-\infty}^0 e^x\, dx\)
\end{question}

%------------------------------------------------------------------------------%
\begin{resolution}{Solução}
Como o intervalo é infinito precisamos explicitar o limite
\[
  \pause
  A
  = \int_{{\color<3>{blue}-\infty}}^0 e^x\, dx
  \pause
  = {\color<3>{blue}\lim_{t\to-\infty}}
    \left(\int_{{\color<3>{blue} t}}^{0} e^x\, dx \;\right)
\]
\end{resolution}

%------------------------------------------------------------------------------%
\begin{resolution}{Primitiva}
  Primeiro vamos calcular a primitiva
  \[
    \pause
    F(x)
    \pause
    \;=\; \int e^x\, dx
    \pause
    \;=\; e^x + c
  \]
\end{resolution}

%------------------------------------------------------------------------------%
\begin{resolution}{Integral definida}
  \begin{align*}
    \onslide<+->{I(t) & = \int_t^0 e^x\, dx  \\[1mm]}
    \onslide<+->{     & = F(x)\evalat{t}{0}  \\[1mm]}
    \onslide<+->{     & = e^x \evalat{t}{0}  \\[1mm]}
    \onslide<+->{     & = e^0 - e^t          \\[1mm]}
    \onslide<+->{     & = 1 - e^t                   }
  \end{align*}
\end{resolution}

%------------------------------------------------------------------------------%
\begin{resolution}{Calculando o limite}
  \begin{align*}
    \onslide<+->{A & = \lim_{t\to-\infty} I(t)                       \\[1mm]}
    \onslide<+->{  & = \lim_{t\to-\infty} \left(1 - e^t \right)      \\[1mm]}
    \onslide<+->{  & = \lim_{t\to-\infty} 1 - \lim_{t\to-\infty}e^t  \\[1mm]}
    \onslide<+->{  & = 1}
  \end{align*}
\end{resolution}

%------------------------------------------------------------------------------%